\documentclass[12pt,oneside]{jlreq}
\pagestyle{empty}
\usepackage[default]{fontsetup}
\usepackage{derivative,fixdif}
\DeclareMathOperator{\atantwo}{atan2}
\newcommand{\bm}[1]{\symbfit{#1}}
\setlength{\paperwidth}{597pt}
\setlength{\paperheight}{845pt}
\setlength{\parindent}{0pt}
\setlength{\hoffset}{-1in}
\setlength{\voffset}{-1in}
\setlength{\oddsidemargin}{0pt}
\setlength{\topmargin}{0pt}
\setlength{\headheight}{0pt}
\setlength{\headsep}{0pt}
\setlength{\textheight}{\paperheight}
\setlength{\textwidth}{\paperwidth}
\begin{document}
値域が通常と異なる関数
\begin{equation*}
    \left\{
        \begin{gathered}
            \arctan(x) : \mathbb{R}\mapsto[0,\pi)\\
            \atantwo(y,x) : \mathbb{R}^2\setminus\{(0,0)\}\mapsto[0,2\pi)
        \end{gathered}
    \right.
\end{equation*}
\section{2次元直交座標}
\begin{gather*}
    \left\{
        \begin{aligned}
            \bm{e}_x&=\left(1,0\right)\\
            \bm{e}_y&=\left(0,1\right)
        \end{aligned}
    \right.
\end{gather*}
\(\bm{e}_x,\bm{e}_y\)は\(t\)に対して不変(定数のように扱われる)\\
\begin{align*}
\bm{r}(x(t),y(t))&=x(t)\bm{e}_x+y(t)\bm{e}_{y}\\
\odv{\bm{r}(x(t),y(t))}{t}&=\odv{x(t)}{t}\bm{e}_x+\odv{y(t)}{t}\bm{e}_{y}\\
\odv[2]{\bm{r}(x(t),y(t))}{t}&=\odv[2]{x(t)}{t}\bm{e}_x+\odv[2]{y(t)}{t}\bm{e}_{y}
\end{align*}
\section{基底ベクトルの変換}
直交座標からの変数変換
\[
(x_1(t),\dots,x_n(t))\to(x_1((y_1(t),\dots,y_n(t))),\dots,x_n((y_1(t),\dots,y_n(t))))
\]
により\[\bm{r}(y_1(t),\dots,y_n(t))=\displaystyle\sum_{k=1}^{n}x_k(y_1(t),\dots,y_n(t))\bm{e}_{x_k}\]として表現可能な場合、各変数方向の基底ベクトルは
\begin{align*}
    \bm{e}_{y_k}(y_1(t),\dots,y_n(t))&=\frac{\pdif{y_k}\bm{r}(y_1(t),\dots,y_n(t))}{\left|\pdif{y_k}\bm{r}(y_1(t),\dots,y_n(t))\right|}\\
    &=\frac{\pdif{y_k}\bm{r}(y_1(t),\dots,y_n(t))}{\sqrt{\displaystyle\sum_{i=1}^{n}\pdif{y_k}{x_i}^2(y_1(t),\dots,y_n(t))}}
\end{align*}
で表される。
\section{2次元極座標}
基本の変換
\begin{gather*}
    \left\{
        \begin{aligned}
            x(r(t),\theta(t))&=r(t)\cos\theta(t)\\
            y(r(t),\theta(t))&=r(t)\sin\theta(t)
        \end{aligned}
    \right.\\
    \left\{
        \begin{aligned}
            r(x(t),y(t))&=\sqrt{x^2(t)+y^2(t)}\\
            \theta(x(t),y(t))&=\atantwo\left(y(t),x(t)\right)
        \end{aligned}
    \right.
\end{gather*}
2次元極座標では
\[
\bm{r}(r(t),\theta(t))=x(r(t),\theta(t))\bm{e}_x+y(r(t),\theta(t))\bm{e}_{y}=r(t)\cos\theta(t)\bm{e}_x+r(t)\sin\theta(t)\bm{e}_y
\]
の関係を用いると、
\begin{align*}
    \bm{e}_{r}(r(t),\theta(t))&=\cos\theta(t)\bm{e}_x+\sin\theta(t)\bm{e}_y\\
    \bm{e}_{\theta}(r(t),\theta(t))&=-\sin\theta(t)\bm{e}_x+\cos\theta(t)\bm{e}_y
\end{align*}
として表される。
次に、各基底ベクトルの時間微分を求める。多変数関数の微分
\[
\odv{f(x_1,\dots,x_n)}{t}=\sum_{i=1}^{n}\pdv{f(x_1,\dots,x_n)}{x_i}\odv{x_i}{t}
\]
を用いると、\(r(t)\)での偏微分は0になり、
\begin{align*}
    \odv{\bm{e}_{r}(r(t),\theta(t))}{t}&=-\sin\theta(t)\odv{\theta(t)}{t}\bm{e}_x+\cos\theta(t)\odv{\theta(t)}{t}\bm{e}_y\\
    &=\odv{\theta(t)}{t}\bm{e}_{\theta}(r(t),\theta(t))\\
    \odv{\bm{e}_{\theta}(r(t),\theta(t))}{t}&=-\cos\theta(t)\odv{\theta(t)}{t}\bm{e}_x-\sin\theta(t)\odv{\theta(t)}{t}\bm{e}_y\\
    &=-\odv{\theta(t)}{t}\bm{e}_{r}(r(t),\theta(t))
\end{align*}
が得られる。この結果を用いて、位置ベクトルの微分を求めることができる。

(積の微分を用いる)
\begin{align*}
    \bm{r}(r(t),\theta(t))&=r(t)\bm{e}_r(r(t),\theta(t))\\
    \odv{\bm{r}(r(t),\theta(t))}{t}&=\odv{r(t)}{t}\bm{e}_r(r(t),\theta(t))+r(t)\odv{\bm{e}_r(r(t),\theta(t))}{t}\\
    &=\odv{r(t)}{t}\bm{e}_r(r(t),\theta(t))+r(t)\odv{\theta(t)}{t}\bm{e}_{\theta}(r(t),\theta(t))\\
    \odv[2]{\bm{r}(r(t),\theta(t))}{t}&=\odv*{\left(\odv{r(t)}{t}\bm{e}_r(r(t),\theta(t))+r(t)\odv{\theta(t)}{t}\bm{e}_{\theta}(r(t),\theta(t))\right)}{t}\\
    &=\odv[2]{r(t)}{t}\bm{e}_r(r(t),\theta(t))+\odv{r(t)}{t}\odv{\bm{e}_r(r(t),\theta(t))}{t}\\
    &+\odv{r(t)}{t}\odv{\theta(t)}{t}\bm{e}_{\theta}(r(t),\theta(t))+r(t)\odv[2]{\theta(t)}{t}\bm{e}_{\theta}(r(t),\theta(t))+r(t)\odv{\theta(t)}{t}\odv{\bm{e}_{\theta}(r(t),\theta(t))}{t}\\
    &=\odv[2]{r(t)}{t}\bm{e}_r(r(t),\theta(t))+\odv{r(t)}{t}\odv{\theta(t)}{t}\bm{e}_{\theta}(r(t),\theta(t))\\
    &+\odv{r(t)}{t}\odv{\theta(t)}{t}\bm{e}_{\theta}(r(t),\theta(t))+r(t)\odv[2]{\theta(t)}{t}\bm{e}_{\theta}(r(t),\theta(t))-r(t)\left(\odv{\theta(t)}{t}\right)^2\bm{e}_{r}(r(t),\theta(t))\\
    &=\left(\odv[2]{r(t)}{t}-r(t)\left(\odv{\theta(t)}{t}\right)^2\right)\bm{e}_r(r(t),\theta(t))\\
    &+\left(2\odv{r(t)}{t}\odv{\theta(t)}{t}+r(t)\odv[2]{\theta(t)}{t}\right)\bm{e}_{\theta}(r(t),\theta(t))
\end{align*}
\section{3次元直交座標}
\begin{gather*}
    \left\{
        \begin{aligned}
            \bm{e}_x&=\left(1,0,0\right)\\
            \bm{e}_y&=\left(0,1,0\right)\\
            \bm{e}_z&=\left(0,0,1\right)
        \end{aligned}
    \right.
\end{gather*}
\(\bm{e}_x,\bm{e}_y,\bm{e}_z\)は\(t\)に対して不変(定数のように扱われる)\\
\begin{align*}
\bm{r}(x(t),y(t),z(t))&=x(t)\bm{e}_x+y(t)\bm{e}_{y}+z(t)\bm{e}_z\\
\odv{\bm{r}(x(t),y(t),z(t))}{t}&=\odv{x(t)}{t}\bm{e}_x+\odv{y(t)}{t}\bm{e}_{y}+\odv{z(t)}{t}\bm{e}_{z}\\
\odv[2]{\bm{r}(x(t),y(t),z(t))}{t}&=\odv[2]{x(t)}{t}\bm{e}_x+\odv[2]{y(t)}{t}\bm{e}_{y}+\odv[2]{z(t)}{t}\bm{e}_{z}
\end{align*}
\section{3次元極座標(右手系)}
基本の変換
\begin{gather*}
    \left\{
        \begin{aligned}
            x(r(t),\theta(t),\varphi(t))&=r(t)\sin\theta(t)\cos\varphi(t)\\
            y(r(t),\theta(t),\varphi(t))&=r(t)\sin\theta(t)\sin\varphi(t)\\
            z(r(t),\theta(t),\varphi(t))&=r(t)\cos\theta(t)
        \end{aligned}
    \right.\\
    \left\{
        \begin{aligned}
            r(x(t),y(t),z(t))&=\sqrt{x^2(t)+y^2(t)+z^2(t)}\\
            \theta(x(t),y(t),z(t))&=\arctan\frac{\sqrt{x^2(t)+y^2(t)}}{z(t)}\\
            \varphi(x(t),y(t),z(t))&=\atantwo\left(y(t),x(t)\right)
        \end{aligned}
    \right.
\end{gather*}
やることは基本一緒なので、結果だけ載せる。
\[
\bm{r}(r(t),\theta(t),\varphi(t))=r(t)\sin\theta(t)\cos\varphi(t)\bm{e}_x+r(t)\sin\theta(t)\sin\varphi(t)\bm{e}_y+r(t)\cos\theta(t)\bm{e}_z
\]
\begin{align*}
    \bm{e}_{r}(r(t),\theta(t),\varphi(t))&=\sin\theta(t)\cos\varphi(t)\bm{e}_x+\sin\theta(t)\sin\varphi(t)\bm{e}_y+\cos\theta(t)\bm{e}_z\\
    \bm{e}_{\theta}(r(t),\theta(t),\varphi(t))&=\cos\theta(t)\cos\varphi(t)\bm{e}_x+\cos\theta(t)\sin\varphi(t)\bm{e}_y-\sin\theta(t)\bm{e}_z\\
    \bm{e}_{\varphi}(r(t),\theta(t),\varphi(t))&=-\sin\varphi(t)\bm{e}_x+\cos\varphi(t)\bm{e}_y
\end{align*}
\begin{align*}
    \odv{\bm{e}_{r}(r(t),\theta(t),\varphi(t))}{t}&=\odv{\theta(t)}{t}\bm{e}_{\theta}(r(t),\theta(t),\varphi(t))+\sin\theta(t)\odv{\varphi(t)}{t}\bm{e}_{\varphi}(r(t),\theta(t),\varphi(t))\\
    \odv{\bm{e}_{\theta}(r(t),\theta(t),\varphi(t))}{t}&=-\odv{\theta(t)}{t}\bm{e}_{r}(r(t),\theta(t),\varphi(t))+\cos\theta(t)\odv{\varphi(t)}{t}\bm{e}_{\varphi}(r(t),\theta(t),\varphi(t))\\
    \odv{\bm{e}_{\varphi}(r(t),\theta(t),\varphi(t))}{t}&=-\sin\theta(t)\odv{\varphi(t)}{t}\bm{e}_{r}(r(t),\theta(t),\varphi(t))-\cos\theta(t)\odv{\varphi(t)}{t}\bm{e}_{\theta}(r(t),\theta(t),\varphi(t))
\end{align*}
\begin{align*}
    \bm{r}(r(t),\theta(t),\varphi(t))&=r(t)\bm{e}_r(r(t),\theta(t),\varphi(t))\\
    \odv{\bm{r}(r(t),\theta(t),\varphi(t))}{t}&=\odv{r(t)}{t}\bm{e}_{r}(r(t),\theta(t),\varphi(t))+r(t)\odv{\theta(t)}{t}\bm{e}_{\theta}(r(t),\theta(t),\varphi(t))+r(t)\sin\theta(t)\odv{\varphi(t)}{t}\bm{e}_{\varphi}(r(t),\theta(t),\varphi(t))\\
    \odv[2]{\bm{r}(r(t),\theta(t),\varphi(t))}{t}&=\left(\odv[2]{r(t)}{t}-r(t)\left(\odv{\theta(t)}{t}\right)^2-r(t)\sin^2\theta(t)\left(\odv{\varphi(t)}{t}\right)^2\right)\bm{e}_r(r(t),\theta(t),\varphi(t))\\
    &+\left(2\odv{r(t)}{t}\odv{\theta(t)}{t}+r(t)\odv[2]{\theta(t)}{t}-r(t)\sin\theta(t)\cos\theta(t)\left(\odv{\varphi(t)}{t}\right)^2\right)\bm{e}_{\theta}(r(t),\theta(t),\varphi(t))\\
    &+\left(2\odv{r(t)}{t}\sin\theta(t)\odv{\varphi(t)}{t}+2r(t)\cos\theta(t)\odv{\theta(t)}{t}\odv{\varphi(t)}{t}+r(t)\sin\theta(t)\odv[2]{\varphi(t)}{t}\right)\bm{e}_{\varphi}(r(t),\theta(t),\varphi(t))
\end{align*}
\section{円筒座標(右手系)}
基本の変換
\begin{gather*}
    \left\{
        \begin{aligned}
            x(\rho(t),\varphi(t),z(t))&=\rho(t)\cos\varphi(t)\\
            y(\rho(t),\varphi(t),z(t))&=\rho(t)\sin\varphi(t)\\
            z(\rho(t),\varphi(t),z(t))&=z(t)
        \end{aligned}
    \right.\\
    \left\{
        \begin{aligned}
            \rho(x(t),y(t),z(t))&=\sqrt{x^2(t)+y^2(t)}\\
            \varphi(x(t),y(t),z(t))&=\atantwo(y(t),x(t))\\
            z(x(t),y(t),z(t))&=z(t)
        \end{aligned}
    \right.
\end{gather*}
円筒座標は2次元極座標と直交座標の和による合成である。
\[
\bm{r}(\rho(t),\varphi(t),z(t))=\rho(t)\cos\varphi(t)\bm{e}_x+\rho(t)\sin\varphi(t)\bm{e}_y+z(t)\bm{e}_z
\]
\begin{align*}
    \bm{e}_{\rho}(\rho(t),\varphi(t),z(t))&=\cos\varphi(t)\bm{e}_x+\sin\varphi(t)\bm{e}_y\\
    \bm{e}_{\varphi}(\rho(t),\varphi(t),z(t))&=-\sin\varphi(t)\bm{e}_x+\cos\varphi(t)\bm{e}_y\\
    \bm{e}_{z}(\rho(t),\varphi(t),z(t))&=\bm{e}_z
\end{align*}
\begin{align*}
    \odv{\bm{e}_{\rho}(\rho(t),\varphi(t),z(t))}{t}&=\odv{\varphi(t)}{t}\bm{e}_{\varphi}(r(t),\theta(t),\varphi(t))\\
    \odv{\bm{e}_{\varphi}(\rho(t),\varphi(t),z(t))}{t}&=-\odv{\varphi(t)}{t}\bm{e}_{\rho}(r(t),\theta(t),\varphi(t))\\
    \odv{\bm{e}_{z}(\rho(t),\varphi(t),z(t))}{t}&=\bm{0}
\end{align*}
\begin{align*}
    \bm{r}(\rho(t),\varphi(t),z(t))&=\rho(t)\bm{e}_{\rho}(r(t),\theta(t),\varphi(t))+z(t)\bm{e}_z\\
    \odv{\bm{r}(\rho(t),\varphi(t),z(t))}{t}&=\odv{\rho(t)}{t}\bm{e}_{\rho}(\rho(t),\varphi(t),z(t))+\rho(t)\odv{\varphi(t)}{t}\bm{e}_{\varphi}(\rho(t),\varphi(t),z(t))+\odv{z(t)}{t}\bm{e}_z\\
    \odv[2]{\bm{r}(\rho(t),\varphi(t),z(t))}{t}&=\left(\odv[2]{\rho(t)}{t}-\rho(t)\left(\odv{\varphi(t)}{t}\right)^2\right)\bm{e}_{\rho}(\rho(t),\varphi(t),z(t))\\
    &+\left(2\odv{\rho(t)}{t}+\odv{\varphi(t)}{t}+\rho(t)\odv[2]{\varphi(t)}{t}\right)\bm{e}_{\varphi}(\rho(t),\varphi(t),z(t))\\
    &+\odv[2]{z(t)}{t}\bm{e}_z
\end{align*}
\end{document}